\chapter{Instalación}

\section{Hardware}

\begin{itemize}
  \item Placa \textbf{PolarFire SoC Discovery Kit} (MPFS095T).
  \item Tarjeta microSD con la imagen de Linux de la placa.
  \item Cable Ethernet a la red de la PC cliente.
  \item Cable USB, sólo para programar el bitstream (JTAG) y para la consola
    serie.
\end{itemize}

Un dispositivo entregado ya trae el bitstream grabado y el software
instalado: basta conectarlo a la red. Las secciones siguientes describen la
instalación desde cero.

\section{Programar el bitstream}

El bitstream se graba por JTAG desde una PC con Microchip Libero SoC o
FlashPro Express, con el programador embebido de la placa conectado por USB.
La FPGA lo conserva sin alimentación.

\begin{lstlisting}[language=shell]
# con FlashPro Express y el archivo de programación del release
FPExpress  # abrir mpeg2fpga-<versión>.job y ejecutar PROGRAM
\end{lstlisting}

Una vez arrancado Linux, el driver informa qué bitstream está cargado:

\begin{lstlisting}[language=shell]
ssh root@placa cat /sys/bus/platform/drivers/mpeg2fpga/*/build
0.1.0+8004bc9
\end{lstlisting}

Es la versión del release y el commit del que se sintetizó; termina en
\cmd{-dirty} si se sintetizó con cambios sin registrar.

\section{Software de la placa}

Primero el driver y su overlay de device tree, después el servicio:

\begin{lstlisting}[language=shell]
scp -r driver/mpeg2fpga/tools root@placa:/tmp/m2f-driver
ssh root@placa sh /tmp/m2f-driver/install-on-board.sh

scp -r daemon api/PROTOCOL-v1.md root@placa:/tmp/m2f-daemon/
ssh root@placa sh /tmp/m2f-daemon/daemon/install-on-board.sh
\end{lstlisting}

El instalador del servicio:
\begin{itemize}
  \item crea el usuario de sistema \cmd{mpeg2fpgad} (sin privilegios);
  \item instala el servicio en \cmd{/usr/local/lib/mpeg2fpgad} y el conversor
    nativo \cmd{libm2fconv.so};
  \item genera el token de acceso en \cmd{/etc/mpeg2fpgad/token};
  \item habilita y arranca la unidad de systemd \cmd{mpeg2fpgad.service}.
\end{itemize}

Opciones: \cmd{-{}-tls} para HTTPS y \cmd{-{}-debug} para los endpoints de
diagnóstico (ver capítulo~\ref{cap:seguridad}).

\section{Verificar}

\begin{lstlisting}[language=shell]
curl http://192.168.18.5:8080/v1/health
{"ok": true}
\end{lstlisting}

\section{Librería cliente}

En la PC, con Python 3.9 o posterior:

\begin{lstlisting}[language=shell]
pip install mpeg2fpga-<versión>-py3-none-any.whl
ssh root@placa cat /etc/mpeg2fpgad/token > token
\end{lstlisting}
